\section{Diagonal isoduality and the logical-Clifford phase}
\label{sec:logical}

The second application uses the atlas gauges to determine the exact
fixed-party transversal logical group of an MDS--CSS encoder.
Throughout this section \(q\) is an odd prime.  Regard any one party of
an \(\AME(6,q)\) stabilizer tensor as an input and the other five as
outputs.  This is a \([[5,1,3]]_q\) encoding isometry.  In prime local
dimension, the symplectic action of a one-qudit Clifford is a block in
\(\mathrm{SL}_2(q)\), so the kernel of the party action on Clifford
tensor symmetries is
\[
 \mathcal K_C=\{(F_1,\ldots,F_6)\in\mathrm{SL}_2(q)^6:
                 (\bigoplus_iF_i)L_C=L_C\}.
\]
For \(q=p^e\) with \(e>1\), the full local Clifford group acts instead
through \(\mathrm{Sp}_{2e}(\F_p)\); the displayed field-linear block
group is then not the full local Clifford kernel.
For an encoder view, the \emph{fixed-party logical symplectic image} means
the image of \(\mathcal K_C\) on its input block.
Corollary~\ref{cor:diagonal-isodual-transversal-group} identifies the
all-length phase boundary as diagonal isoduality.  For six-arcs, classical
self-association turns that intrinsic code condition into the conic
criterion below.

\begin{lemma}[Fixed-label self-association]\label{lem:six-arc-self-association}
Let \(H\) be a rank-three \(3\times6\) matrix whose columns are a
six-arc, and put \(C=\ker H\).  There is a nonsingular diagonal matrix
\(S\) with \(SC=C^\perp\) if and only if the six columns of \(H\) lie on
a conic.
\end{lemma}

\begin{proof}
Let \(G\) be a \(3\times6\) generator matrix for \(C\).  Its columns are
the labeled Gale transform of the columns of \(H\).  Since
\(C^\perp=\operatorname{row}(H)\),
\[
 SC=C^\perp
 \quad\Longleftrightarrow\quad
 \operatorname{row}(GS)=\operatorname{row}(H)
 \quad\Longleftrightarrow\quad
 AGS=H\ \text{ for some }A\in\mathrm{GL}_3(q).
\]
Thus the condition is precisely fixed-label projective equivalence to the
Gale transform.  This fixed locus can be identified directly over the
finite field.  Put
\[
 V_i=(x_i^2,y_i^2,z_i^2,x_iy_i,x_iz_i,y_iz_i)
\]
and let \(V\) be the \(6\)-by-\(6\) matrix with rows \(V_i\).
Its determinant vanishes exactly when a nonzero quadratic form contains
all six columns.  Every five rows of \(V\) are independent: a relation
supported on at most five points would give
\(H_I D H_I^T=0\) for a nonsingular diagonal \(D\) on its support \(I\);
for \(|I|\geq3\), this would put the three-dimensional row space of
\(H_I D\) inside \(\ker H_I\), which has dimension at most two, while
smaller supports are immediate.  Hence any relation
\(\sum_i d_i V_i=0\) has every \(d_i\ne0\).

Writing \(D=\diag(d_i)\), that relation is \(HDH^T=0\).  Therefore
\(\operatorname{row}(HD)\subseteq\ker H=C\), and equality follows from
dimension three.  Thus \(CD^{-1}=C^\perp\).  Conversely, a nonsingular
diagonal equivalence \(SC=C^\perp\) gives
\(HS^{-1}H^T=0\), hence a relation among the rows of \(V\) and a conic
through the six columns.  The conic is irreducible because the columns
form an arc.  This is the finite-field fixed-label form of the classical
six-point self-association theorem
\cite[Section~2.3]{HowardMillsonSnowdenVakil2008}.
\end{proof}

\begin{theorem}[Fixed-party logical phase]\label{thm:logical-phase}
For an equal-phase CSS state from a six-arc over an odd prime field, the
fixed-party logical symplectic image is
\[
\begin{cases}
\mathrm{SL}_2(q),&\text{if the arc is conic/GRS},\\
T=\{\diag(a,a^{-1}):a\in\F_q^\times\},&\text{if the arc is nonconic}.
\end{cases}                                             \tag{5.1}
\]
Thus every encoder view of a GRS tensor realizes the full one-qudit
logical symplectic group from fixed-party symmetries.  Off the GRS locus,
the fixed-party logical symplectic image is the split torus.  At a chosen
anchor leg, any party-moving isoduality that exchanges the CSS \(X\)- and
\(Z\)-spaces transports the encoder view by the nontrivial Weyl-group
element, so the possible anchor blocks generate the normalizer of \(T\)
together with \(T\).  This is an action on one fixed encoder only when
the isoduality fixes its input leg.
\end{theorem}

\begin{proof}
Corollary~\ref{cor:diagonal-isodual-transversal-group} gives
\(\mathrm{SL}_2(q)\) exactly on the diagonally isodual branch and \(T\)
otherwise.  Lemma~\ref{lem:six-arc-self-association} identifies diagonal
isoduality with the conic locus.  At a chosen anchor leg, a party-moving
isoduality exchanges the CSS \(X\)- and \(Z\)-axes and hence represents
the nontrivial Weyl-group element normalizing \(T\).  It acts on one fixed
encoder only when the party motion fixes that encoder's input.
\end{proof}

Adding logical Paulis recovers the exact fixed-party projective logical
groups in Corollary~\ref{cor:diagonal-isodual-transversal-group}.

This is an operational distinction, not a classical automorphism count.
The fixed-party non-GRS group has order \(q^2(q-1)\), while the GRS
group has order \(q^2\lvert\mathrm{SL}_2(q)\rvert\).

The MDS property also makes minimum operator pushing uniform.  For any
input party and nonidentity Pauli label, each of the ten four-sets
containing that party gives a unique stabilizer representative supported
there.  Removing the input leaves a weight-three output Pauli.  Thus each
input Pauli has exactly ten minimum-support pushes, although the symmetry
orbits of those pushes can distinguish finer projective classes.  This
uses only the projection isomorphism for the shortened planes, not the
finite group calculation.
